\section{Conclusion}
\label{sec:conclusion}

\starname{} is a failure-aware state-conditioned routing formulation for specialist composition.
Its central argument is that compositional reasoning over heterogeneous specialists requires a routing policy conditioned on \emph{what went wrong}: \sfail, \sblock, and \smiss demand semantically distinct recovery, and success-only routing is structurally unable to learn any of them (Theorem~\ref{thm:recovery}).
On three spatiotemporal benchmarks and eight backbones, the routing formulation consistently improves per-benchmark EM over LLM-only baselines and outperforms Reflexion, ReAct, and other prompting and agent baselines.
Router-specific ablations isolate the contribution of each routing component and a dedicated recovery analysis provides the sharpest evidence: STARK queries entering the $\smiss$ state recover at 80.8\%, above the no-failure baseline, suggesting that recovery transitions rescue mis-routed queries rather than provide generic fallback.
Future work targets the parameter-extraction step that bounds \sfail{} recovery and learned refinements of the routing policy beyond trace-count estimation.

\paragraph{Limitations.}
STAR currently learns routing matrices from execution traces within a given task taxonomy.
As the transfer experiments show, routing generalizes best when source and target benchmarks share computational primitives and specialist dependencies.
Future work will study task-agnostic recovery abstractions and larger composite-query evaluations.
